\section{Úvodem}\label{sec:kompilatory-uvod}

\emph{Kompilátor} je program,~který překládá program zapsaný ve zdrojovém jazyce do ekvivalentního programu v~cílovém jazyce. Vstup obvykle nazýváme \emph{zdrojovým kódem} a~výstup \emph{cílovým kódem}.

Cílovým jazykem nemusí být přímo strojový kód konkrétního procesoru. Překladač může vytvářet assembler,~bajtkód pro virtuální stroj,~jiný programovací jazyk nebo obecnější mezikód. Podstatné je,~aby se zachoval význam programu: pro stejné vstupy má cílový program poskytovat stejné výsledky jako program zdrojový.

\begin{figure}[h]
    \centering
    \begin{tikzpicture}[node distance=24mm]
    \node[task,fill=blue!10,minimum width=31mm] (source) {zdrojový\\kód};
    \node[task,fill=orange!12,minimum width=31mm,right=of source] (compiler) {kompilátor};
    \node[task,fill=green!10,minimum width=31mm,right=of compiler] (target) {cílový\\kód};
    \draw[diredge] (source)--(compiler);
    \draw[diredge] (compiler)--(target);
    \node[below=4mm of source,align=center,font=\small]
        {\texttt{C}, \texttt{Java},\\\texttt{C\#}, \dots};
    \node[below=4mm of target,align=center,font=\small]
        {assembler, mezikód,\\strojový kód};
\end{tikzpicture}

    \caption{Základní pohled na kompilátor.}
    \label{fig:kompilator-zaklad}
\end{figure}

\subsection{Překladový řetězec}

Slovem \emph{kompilátor} někdy neformálně označujeme celý systém,~který vytvoří spustitelný program. Ve skutečnosti může být práce rozdělena mezi několik nástrojů. Typický překlad jednoduchého programu v~jazyce C znázorňuje obrázek \ref{fig:prekladovy-retezec-c}.

\begin{figure}[h]
    \centering
    \begin{tikzpicture}[
    file/.style={task,fill=blue!8,minimum width=22mm,font=\small},
    tool/.style={task,fill=orange!12,minimum width=25mm,
        text width=21mm,align=center,font=\small}
]
    \node[file] (c) at (0,0) {\texttt{hello.c}};
    \node[tool] (pre) at (3,0) {preprocesor};
    \node[file] (i) at (6,0) {\texttt{hello.i}};
    \node[tool] (comp) at (9,0) {kompilátor};
    \node[file] (s) at (12,0) {\texttt{hello.s}};

    \node[tool] (asm) at (9,-1.8) {assembler};
    \node[file] (o) at (6,-1.8) {\texttt{hello.o}};
    \node[tool] (link) at (3,-1.8) {linker};
    \node[file,fill=green!10] (exe) at (0,-1.8) {\texttt{hello}};

    \draw[diredge] (c)--(pre);
    \draw[diredge] (pre)--(i);
    \draw[diredge] (i)--(comp);
    \draw[diredge] (comp)--(s);
    \draw[diredge] (s.south)--++(0,-0.45)-|(asm.east);
    \draw[diredge] (asm)--(o);
    \draw[diredge] (o)--(link);
    \draw[diredge] (link)--(exe);
\end{tikzpicture}

    \caption{Zjednodušený překladový řetězec jazyka C.}
    \label{fig:prekladovy-retezec-c}
\end{figure}

Jednotlivé části mají následující úlohy:
\begin{itemize}
    \item \emph{preprocesor} zpracuje direktivy,~například vložení hlavičkového souboru pomocí \texttt{\#include};
    \item vlastní \emph{kompilátor} převede předzpracovaný program do assembleru;
    \item \emph{assembler} vytvoří objektový soubor se strojovými instrukcemi a~pomocnými informacemi;
    \item \emph{linker} neboli sestavovací program spojí objektové soubory a~potřebné knihovny do výsledného programu.
\end{itemize}

Uživatel obvykle spustí jediný ovladač překladače,~například \texttt{gcc},~který jednotlivé nástroje zavolá za něj. Rozdělení je však důležité při hledání chyb: chyba kompilace,~chybějící symbol při linkování a~chyba vzniklá až za běhu programu jsou tři různé situace.

\subsection{Strojový kód a~mezikód}

\emph{Strojový kód} obsahuje instrukce konkrétního procesoru. Procesor jej může vykonávat přímo,~ale stejný soubor zpravidla nelze přenést na počítač s~jinou instrukční sadou nebo jiným operačním systémem.

\emph{Mezikód} je vnitřní reprezentace programu ležící mezi zdrojovým a~strojovým kódem. Může být navržen hlavně pro další analýzu a~optimalizaci nebo jako přenositelný bajtkód pro virtuální stroj. Díky mezikódu nemusíme pro každou dvojici zdrojového jazyka a~cílové platformy vytvářet celý překladač znovu: přední část přeloží různé jazyky do společné reprezentace a~koncová část ji převede pro konkrétní počítač.

\subsection{Fáze překladače}

Překladač zdrojový program obvykle nezpracuje jediným obřím krokem. Postupně jej převádí do reprezentací,~se kterými se stále snáze pracuje.

\begin{figure}[h]
    \centering
    \begin{tikzpicture}[
    phase/.style={task,minimum width=27mm,font=\small},
    data/.style={task,minimum width=24mm,font=\small}
]
    \draw[thick,rounded corners=4pt] (-1.5,1.45) rectangle (10.7,-2.75);
    \node[fill=white,inner xsep=6pt] at (4.6,1.45) {\textbf{překladač}};

    \node[data,fill=red!12] (src) at (0,0) {zdrojový\\kód};
    \node[phase,fill=red!12] (lex) at (3,0) {lexikální\\analýza};
    \node[phase,fill=red!12] (syn) at (6,0) {syntaktická\\analýza};
    \node[phase,fill=red!12] (sem) at (9,0) {sémantická\\analýza};

    \node[data,fill=orange!15] (ir) at (9,-1.6) {mezikód};
    \node[phase,fill=orange!15] (opt) at (6,-1.6) {optimalizace};
    \node[phase,fill=orange!15] (gen) at (3,-1.6) {generování\\kódu};
    \node[data,fill=green!12] (target) at (0,-1.6) {cílový\\kód};

    \draw[diredge] (src)--(lex);
    \draw[diredge] (lex)--(syn);
    \draw[diredge] (syn)--(sem);
    \draw[diredge] (sem)--(ir);
    \draw[diredge] (ir)--(opt);
    \draw[diredge] (opt)--(gen);
    \draw[diredge] (gen)--(target);

    \node[font=\small] at (4.5,0.85) {\textbf{přední část (front end)}};
    \node[font=\small] at (4.5,-2.35) {\textbf{koncová část (back end)}};
\end{tikzpicture}

    \caption{Základní fáze překladače.}
    \label{fig:faze-prekladace}
\end{figure}

\begin{enumerate}
    \item \emph{Lexikální analýza} seskupí znaky do tokenů.
    \item \emph{Syntaktická analýza} z~tokenů sestaví strom podle gramatiky.
    \item \emph{Sémantická analýza} zkontroluje například deklarace,~typy a~platnost operací.
    \item Z~ověřeného stromu vznikne \emph{mezikód}.
    \item \emph{Optimalizace} mezikód upraví,~aniž by změnila význam programu.
    \item \emph{Generování kódu} vytvoří cílový kód.
\end{enumerate}

První tři fáze tvoří \emph{přední část překladače} (anglicky \emph{front end}). Ta závisí hlavně na zdrojovém jazyce. Optimalizace a~generování cílového kódu se obvykle řadí do \emph{koncové části} (anglicky \emph{back end}),~která více závisí na cílové platformě. Hranice ale není u~všech překladačů úplně stejná.

\notebox{Jednotlivé fáze nemusí odpovídat samostatným programům ani jedinému průchodu zdrojovým textem. Jde především o~logické rozdělení odpovědností. Praktický překladač může několik fází spojit nebo některou z~nich provést vícekrát.}

\subsection{Průběžný příklad}

V~dalších sekcích vytvoříme malý jazyk s~číselnými výrazy,~deklarací proměnné příkazem \texttt{let} a~výpisem příkazem \texttt{print}. Program
\begin{verbatim}
let x = 2 + 3 * 4;
print x;
\end{verbatim}
má vypsat číslo $14$. Překladač jej převede do jednoduchého zásobníkového bajtkódu. Ten potom vykoná náš virtuální stroj. Příklad je záměrně malý,~ale projde všemi hlavními fázemi z~obrázku \ref{fig:faze-prekladace}.
